% Generated draft: requires mathematical and source review.
\documentclass[11pt]{article}
\usepackage[margin=25mm]{geometry}
\usepackage{fontspec}
\setmainfont{Latin Modern Roman}
\usepackage{amsmath,amssymb,mathtools,unicode-math}
\setmathfont{Latin Modern Math}
\usepackage{enumitem,needspace,xurl,hyperref}
\setlength{\parindent}{0pt}
\setlength{\parskip}{5pt}
\setlength{\emergencystretch}{3em}
\newcommand{\N}{\mathbb{N}}
\newcommand{\Z}{\mathbb{Z}}
\newcommand{\Q}{\mathbb{Q}}
\newcommand{\R}{\mathbb{R}}
\newcommand{\C}{\mathbb{C}}
\newcommand{\F}{\mathbb{F}}
\newcommand{\A}{\mathbb{A}}
\newcommand{\PP}{\mathbb{P}}
\newcommand{\E}{\mathbb{E}}
\newcommand{\eps}{\varepsilon}
\newcommand{\dd}{\,\mathrm{d}}
\begin{document}
\section*{UnsolvedMath ID: 20000419}
\textbf{Source problem code:} catalogue entry
\subsection{Definitions and mathematical statement}
J. Ellenberg's AIM question asks: when rational Tate classes are known to be algebraic, what is the integral cycle-class image, and for which primes \(\ell\) can surjectivity fail? In codimension notation the comparison is between
\[
\operatorname{CH}^i(X)\otimes\Q_\ell\longrightarrow H^{2i}(X_{\overline k},\Q_\ell(i))
\quad\text{and}\quad
\operatorname{CH}^i(X)\otimes\Z_\ell\longrightarrow H^{2i}(X_{\overline k},\Z_\ell(i)),
\]
with the intended Galois-invariant Tate targets and field of definition specified. The original loses superscripts/subscripts in its \(C_i\) notation. Its hypersurface example concerns one-dimensional cycles mapping to \(H^4\) on a threefold, hence codimension two, not divisors. Integral cohomology with torsion and its quotient modulo torsion are also distinct targets.

The original comments attribute to Koll\'ar the statement that for a very general hypersurface \(X\subset\PP^4\) with \(d^3\mid\deg X\), the one-cycle image lies in \(dH^4(X,\Z)\), and say it was not known whether this occurs over \(\overline\Q\). The hypotheses on \(d\) and the precise theorem are not supplied; the unknown-over-\(\overline\Q\) statement is historical source context. C. Schoen's comment says the rational Tate conjecture implies the integral version for one-cycles over finite fields at \(\ell\ne p\); T. Geisser's says a rational-to-integral implication holds for Weil-\'etale cohomology, relating the question to comparison with \'etale cohomology. Their exact scopes and global versus fixed-variety hypotheses are not included and must not be strengthened silently.

The attached report claims that when the finite-field variety is allowed to vary, failure occurs at every prime \(\ell\): projective approximations in Kameko's codimension-two examples have rational degree-four cohomology spanned by algebraic Chern classes but integral cycle-class failure even modulo torsion. No explicit varieties, fields, or full references are supplied.

Its characteristic-zero candidate applies, after a finite field extension, to a very general principally polarized abelian \(g\)-fold over a finitely generated field. With polarization class \(\Theta\), the intermediate minimal-class lines are
\[
\Z_\ell\,\frac{[\Theta]^c}{c!},\qquad 2\leq c\leq g-1
\]
in the usual nondivisor, nontop interpretation. Combining the March 2026 Engel--de Gaay Fortman--Schreieder results with Milne's absolute-Hodge descent, it asserts a \(2\)-adic integral Tate failure on every such line for \(g\geq4\), exact line index \(2\) for \(g=4,5\), and a \(3\)-adic failure on the minimal curve line \(c=g-1\) for \(g\geq7\). The exact intermediate range and line-index normalization must be checked against the cited theorem. A very general finitely generated-field assertion is not an assertion about every abelian variety or automatically about one defined over a number field or \(\overline\Q\).

The report also claims a two-sided determinant bound on integral defects and a conditional finite exceptional-prime set from uniform perfect-pairing data. No determinant formula or uniformity hypotheses are printed. Thus these are additional conditional audit claims, not a theorem that one finite prime set works for all varieties. Rational algebraicity does not itself establish integral saturation.
\subsection{Short English statement}
Investigate which primes obstruct integral saturation of a rationally algebraic Tate lattice. The report cites varying finite-field failures and proposes prime-specific minimal-class failures on very general abelian varieties, while retaining torsion and field-of-definition distinctions.
\subsection{Sources}
\begin{itemize}
\item AIM workshop ``The Tate conjecture,'' Problem 8, attributed to J. Ellenberg: \url{https://aimath.org/WWN/tateconjecture/tateconjecture.pdf}. The verification says OCR lost indices, accents, and \(\ell\); no PDF was accessed here.
\item Local canonical registry \texttt{aim-algebraic-number-theory-notes.json}, \texttt{notes[70]}, tag \texttt{problem}; permanent UnsolvedMath ID \(20000419\), source code \texttt{catalogue entry}. Nested comments attribute the hypersurface example to Koll\'ar, the finite-field one-cycle assertion to C. Schoen, and the Weil-\'etale comparison to T. Geisser. Original remarks and literature field are otherwise empty.
\item Engel--de Gaay Fortman--Schreieder, arXiv:2507.15704v3, \url{https://arxiv.org/abs/2507.15704v3}, cited in the report as March 2026 results; the summary also invokes Milne's absolute-Hodge descent and Kameko's examples without full bibliographic entries.
\item Supplied research attempt 2 proposing the prime-specific corollary and determinant audit. It is machine generated, not peer reviewed, and requires independent verification before citation.
\end{itemize}
\end{document}
